% =========================== 9. SENSITIVITY ===========================
\section{Sensitivity and the Boundaries of Validity}\label{sec:sens}

Four quantities are not identified by the data --- the sampling rate, the initial
inventory, the incentive parameters and the cost parameters --- and one is estimated
with sampling error, the fleet size. We perturb each and report how far the
conclusions move.

\subsection{The fleet estimate}

Applying \eqnref{eq:fleet} day by day gives the estimates below.

\begin{center}
\small
\begin{tabular}{@{}lrrrrr@{}}
\toprule
\textbf{Date (May 2017)} & 10 & 11 & 12 & 15 & 16 \\
\textbf{Fleet estimate} & 313{,}181 & 332{,}398 & 325{,}991 & 336{,}788 & 343{,}605 \\
\midrule
\textbf{Date (May 2017)} & 13 & 14 & 18 & 19 & 21 \\
\textbf{Fleet estimate} & 326{,}046 & 326{,}690 & 292{,}831 & 300{,}887 & 306{,}947 \\
\bottomrule
\end{tabular}
\end{center}

The estimate ranges from 292{,}831 to 343{,}605, with mean 320{,}536 and standard
deviation 15{,}580 --- a coefficient of variation of 4.9\%. The low values on 18 and 19
May coincide with the highest trips-per-bicycle ratios of the sample (1.65 against
1.46), which is what a day of unusually many short trips would produce; the estimate
is sensitive to the tail of the rides-per-bicycle distribution, and a promotion that
adds many single-trip users biases it downward. We adopt the weekday value 343{,}638
and treat fleet level as a decision variable $\alpha$ rather than a fixed fact, which
is what Q2's right-sizing analysis requires anyway.

The unknown sampling rate does not matter. Resampling the orders of 16 May at rates
from 5\% to 50\% and re-estimating gives 330{,}586 to 338{,}401 bicycles, within 4\% of
the full-sample value, exactly as the invariance argument of \S\ref{sec:data} predicts.

\subsection{The incentive parameters}

The three logit parameters of \eqnref{eq:logit} were varied one at a time around the
baseline ($\kappa=2.0$, $\rho=0.35$, $c_0=0.8$), re-solving \eqnref{eq:mix} each time.

\begin{table}[htbp]
\centering
\caption[Sensitivity of the incentive design]{Sensitivity of the incentive design. The all-truck baseline is 752{,}404 CNY
per day.}
\label{tab:q3sens}
\small
\begin{tabular}{@{}lrrrr@{}}
\toprule
\textbf{Variant} & \textbf{Optimal reward} & \textbf{Placements} &
\textbf{Total cost} & \textbf{Saving vs truck} \\
\midrule
Baseline & 1.20 & 115{,}385 & 489{,}578 & 34.9\% \\
Value of money $\kappa=1.0$ & 1.20 & 78{,}709 & 573{,}116 & 23.8\% \\
$\kappa=3.0$ & 1.05 & 131{,}671 & 432{,}730 & 42.5\% \\
Cost of time $\rho=0.20$ & 1.00 & 123{,}685 & 445{,}935 & 40.7\% \\
$\rho=0.50$ & 1.30 & 102{,}666 & 528{,}814 & 29.7\% \\
Inertia $c_0=0.4$ & 1.05 & 117{,}884 & 466{,}202 & 38.0\% \\
$c_0=1.2$ & 1.35 & 112{,}779 & 512{,}429 & 31.9\% \\
\bottomrule
\end{tabular}
\end{table}

The qualitative conclusion is robust: in every variant incentives are worth using, the
optimal reward lies between 1.00 and 1.35 CNY, and the saving over an all-truck policy
lies between 24\% and 43\%. The level of the reward and the size of the programme are
not robust in the same way --- they move with the behavioural parameters by up to a
factor of 1.7 in placements --- which is why we present the design as a function of a
parameter that a pilot study should estimate before the scheme is scaled.

\subsection{The price of a truck-kilometre}

Re-solving the Q2 program at $\alpha=1$, $K=60$ with the truck-kilometre price at 3, 6
and 12 CNY gives:

\begin{center}
\small
\begin{tabular}{@{}lrrr@{}}
\toprule
$c_{\text{km}}$ (CNY/km) & \textbf{Bicycles moved} & \textbf{Truck-km} & \textbf{Cost (CNY/day)} \\
\midrule
3  & 25{,}587 & 5{,}814 & 66{,}242 \\
6  & 23{,}959 & 5{,}441 & 83{,}325 \\
12 & 21{,}835 & 4{,}906 & 113{,}964 \\
\bottomrule
\end{tabular}
\end{center}

A four-fold increase in the truck price reduces the relocation volume by only 15\%,
because the program is constrained by the need to restore balance, not by the budget.
The plan is therefore robust to the most uncertain of the operating costs: getting the
price wrong by a factor of two changes the schedule by a few thousand bicycles, not by
its structure.

\subsection{The initial inventory}

This is the largest single uncertainty, and it is structural rather than parametric.
The data record where each bicycle ended its last observed trip, not where it stood at
midnight. Two estimators are available --- the last-trip-end distribution and the
first-trip-start distribution --- and they agree at the zone level but not at the cell
level, where 5,263 cells are served by 343,638 bicycles and many cells have no evidence
of any bicycle at all.

The consequence is a wide range in the simulated service level at cell resolution. In
the zone-level simulation of \S\ref{sec:q2} no departure is ever unmet; in the
cell-level simulation with a 300 m walking radius, 13.4\% of departures go unserved at
$\alpha=1$, rising to 40.3\% at $\alpha=0.2$. The gap is not a modelling failure but an
identification result: the cell-level figure is dominated by cells whose midnight stock
is unknown, and it is best read as an upper bound on unmet demand. The zone-level
figure, where stocks are large enough to be inferred, is the reliable one, and the LP's
deterministic service metric lies between them.

We use this spread explicitly rather than choosing a point. All statements about
service in this paper are made at the zone level or in the LP; no conclusion depends on
the cell-level figure, which we report because suppressing it would overstate the
precision of the model.

\subsection{Range of validity}

The conclusions hold for a dense, flat, monocentric city of about 1{,}200 km$^2$ with
6.2 rides per bicycle per day, in mild weather, on a weekday. They should not be
extrapolated to:

\begin{itemize}[leftmargin=1.4em]
\item \emph{systems with stations}, where the docking constraint changes the inventory
  process and the rebalancing problem becomes a routing problem with hard capacity;
\item \emph{low-density cities}, where the walking radius of Q4 covers a much smaller
  share of arrivals and the incentive pool of Q3 vanishes;
\item \emph{extreme weather or holidays}, not represented in the ten usable days;
\item \emph{the long run}, where the attraction $A_j$ itself changes with land use. The
  model is a description of an equilibrium, not a forecast of one.
\end{itemize}
